\honest{Decoherence is Simple}{Eliezer Yudkowsky}{2008}

This is an epistle to the physicists. My father, a physicist, told me that physics cannot be understood without its mathematics; I believed Feynman instead, and later found that my father was right. I am a Bayesian, not a physicist, but I must correct the misuse of three words: simple, falsifiable and testable. Here I take the first. What follows is a restatement of ``Belief in the Implied Invisible,'' applied to quantum physics.

The remark I answer, in my paraphrase, is that many-worlds posits vast numbers of other worlds, and Occam's razor says not to multiply entities. \nb{The objection is real; the Stanford Encyclopedia's entry on many-worlds reports it as the source of ``much of the opposition'' to the view.} Those who say this usually add that some apply the razor to laws, not objects. I quote the razor in Latin as ``The original formulation of William of Ockham.'' \nb{As noted for ``Belief in the Implied Invisible,'' the phrasing comes from John Punch in 1639, not Ockham.}

To calculate simplicity there are two roads. One asks how big a program must be to compute a hypothesis's predictions: Kolmogorov complexity and Solomonoff induction. The other asks how far the hypothesis shortens a message describing the data: Minimum Message Length. ``The woman down the street is a witch'' is eleven words, but on either road it buys nothing. You cannot pick and choose between the two, because they ``were proven equivalent.'' \nb{In ``Occam's Razor'' Yudkowsky had called them ``nearly equivalent'' and named the difference. Wallace and Dowe (1999) ``attempt to establish a parallel'' between Minimum Message Length and a two-part version of Kolmogorov's model.} Both count the length of the code, not the memory it uses.

Why should you accept these formalisms? I claim experimental evidence, and begin with the conjunction rule. Each detail a story states splits its probability, so the razor's ``entities'' are the details a theory must state. ``I will win the lottery'' mentions one person and is far less likely than ``someone will win.'' Objects count only when they must be specified one by one, as with particular billiard balls; the exact state of the air in a warm room need not be specified. A gram of hydrogen holds six hundred billion trillion atoms, and we believe it; a theory with that many separate laws we would not.

Why not count memory or running time? First, because resources ``aren't mutually exclusive possibilities.'' \nb{Schmidhuber's speed prior (2002) counts running time and is still a normalized prior; Schmidhuber concluded from it that large-scale quantum computation ``will not work well.''} Second, because history shows it does not work: the razor was raised against the idea that nebulae were galaxies, and the universe keeps getting bigger. \nb{In ``Belief in the Implied Invisible'' Yudkowsky offered the nebulae story from memory, after failing to find a source; the editors could not find one either. The other part of the promised evidence is a remark about researchers' programs that names no study.}

If the expansion of the universe continues, a spaceship will one day pass beyond our cosmological horizon. It does not then vanish. If the razor counted objects, the vanishing model would win; but making the ship vanish needs a new law. \nb{Continued expansion alone would not create the horizon; acceleration does, and current cosmology holds that the expansion accelerates.}

Decoherence, ``a.k.a. many-worlds,'' says measurements follow the same rules as everything else. \nb{Physicists use ``decoherence'' for a process that every interpretation accepts; the Stanford Encyclopedia's entry on it says that decoherence ``as such does not provide a solution to the measurement problem.''} Many questions remain, such as why binary measurements are not all 50/50, but I set them aside. \nb{That question bears on simplicity. Vaidman's entry says the probability rule in many-worlds ``involves adding a postulate,'' and also reports the view that many-worlds is ``at no disadvantage'' here. The next day Yudkowsky described the decoherent theory as the wavefunction's equation ``plus the Born probability rule.''} Many-worlds generates its worlds from compact laws, and simulating any version of quantum mechanics is exponentially costly, many-worlds only ``more so.'' \nb{Hemmo and Shenker argue that the many-worlds postulates alone do not explain our experience, so something must be added. In program-length terms, Hutter (2010) argued that a program must also locate the observer.}

Perhaps critics mistake their shock at a large universe for a probability penalty. The objection ``is entirely defective as probability theory. It is not fixable. It is bad math.'' \nb{The post shows that the objection conflicts with the two formalisms and with history. Probability theory does not choose the prior, so this is a dispute about priors, one in which many philosophers share Yudkowsky's side.}
